% X_0^6(649)* is hyperelliptic: three degree-two CM divisors whose differences are torsion (dual test), with the
% 3-part of the torsion killed by the component group at 3 (Molina's reduction of CM points), so two of them are
% linearly equivalent.  Tests: tests/Hyperelliptic6_649.m and tests/_offline/Hyperelliptic6_649.m.
\documentclass[11pt,reqno]{amsart}
\usepackage[margin=1.15in]{geometry}
\usepackage{amsmath,amssymb,amsthm,mathtools}
\usepackage[colorlinks=true,linkcolor=blue,citecolor=blue,urlcolor=blue]{hyperref}
\theoremstyle{plain}
\newtheorem{theorem}{Theorem}[section]
\newtheorem{lemma}[theorem]{Lemma}
\theoremstyle{definition}
\newtheorem{remark}[theorem]{Remark}
\DeclareMathOperator{\Jac}{Jac}
\DeclareMathOperator{\Pic}{Pic}
\newcommand{\Q}{\mathbb{Q}}
\newcommand{\Z}{\mathbb{Z}}
\newcommand{\F}{\mathbb{F}}
\begin{document}
\title{$X_0^{6}(649)^*$ is hyperelliptic}
\author{}
\date{\today}
\maketitle

\section{Statement}

Let $Y=X_0^6(649)^*$ be the quotient of $X_0^6(649)$ by its full Atkin--Lehner group $W$, of order $16$. It has
genus three and no further quotient, so it is hyperelliptic if and only if it carries two distinct linearly
equivalent divisors of degree two. Consider the three divisors
\[
D_1=P_{-24}+P_{-132},\qquad D_2=P_{-40}+P_{-187},\qquad D_3=P_{-52}+P_{-660},
\]
each the sum of two rational CM points of $Y$.

\begin{theorem}\label{thm:main}
Two of $D_1,D_2,D_3$ are linearly equivalent. In particular $Y$ is hyperelliptic.
\end{theorem}

The three divisors are the fibres of the hyperelliptic map found by the Shimizu screen of the invariant
differentials of $Y$ at its rational CM points; the proof uses them only as divisors.

\section{Proof}

\emph{The differences are torsion.} The dual test of \texttt{paper/hyperelliptic-15-38.tex}, Proposition~2.1,
at $M=4\cdot1947=7788$ (\texttt{tests/\_offline/Hyperelliptic6\_649.m}): a principal part with divisor
$D_i-D_j$ in the dictionary of Guo--Yang pairs to zero with the space $W$ of cusp forms of weight $\tfrac32$ on
$\Gamma_0(7788)$ vanishing at every index carrying no CM point, for each of the three pairs $i<j$; the test
also checks that $T(49)$ acts on $W$ with a factor of the polynomial of the traces of Frobenius of $Y$ at $7$,
and refuses $P_{-24}-P_{-132}$ and $P_{-24}+P_{-40}-P_{-132}-P_{-187}$.

\emph{The torsion of $\Jac(Y)(\Q)$ has order dividing $6$.} The orders $\#\Jac(Y)(\F_p)=L_p(1)$, from the
point counts of $Y$ over $\F_p$, $\F_{p^2}$, $\F_{p^3}$ by the trace formula, have greatest common divisor $6$
over the good primes below $120$.

\emph{The $3$-part injects into $\Phi_3$.} $Y$ is a quotient of the \v{C}erednik--Drinfeld model of $X_0^6(649)$
at $3$, so its stable reduction at $3$ is a union of rational curves meeting at ordinary double points, with
dual graph the quotient by $W$ of the \v{C}erednik--Drinfeld graph; its first Betti number is $3$, the genus,
so the reduction is totally degenerate and the identity component of the N\'eron model of $\Jac(Y)$ at $3$ is
a torus $T$. The torus is split: Frobenius acts on the \v{C}erednik--Drinfeld graph as $w_3$, which lies in
$W$, so it acts trivially on the quotient graph and on $H_1$ of it, the character group of $T$. Hence
$T(\F_3)=(\F_3^\times)^3$ has no $3$-torsion. The kernel of reduction $\Jac(Y)(\Q_3)\to\mathcal J(\F_3)$ is
the group of $3\Z_3$-points of the formal group of $\mathcal J$, which is torsion-free because the formal
logarithm is an isomorphism onto $(3\Z_3)^3$ (the ramification index $1$ is less than $3-1$). So
$\Jac(Y)(\Q)[3^\infty]$ injects into $\mathcal J(\F_3)$ and, having no image in $T(\F_3)$, into the component
group $\Phi_3$, which is the critical group of the quotient graph: $\Phi_3\cong\Z/102$.

\emph{The differences reduce to zero.} Each $P_d$ reduces at $3$ to the singular point or component of the
\v{C}erednik--Drinfeld fibre whose Eichler order admits the optimal embedding of the order of discriminant $d$
(Molina, \emph{Ribet bimodules and the specialization of Heegner points}, Theorems 5.3, 5.4, 5.8): to a singular
point when $3$ ramifies in $\Q(\sqrt d)$ ($d=-24,-132,-660$), each at a singular point whose branches are
swapped by an involution of $W$ divisible by $3$, hence at a smooth point of a merged component of $Y$; to a
component when $3$ is inert ($d=-40,-187,-52$). The classes of $D_1-D_2$, $D_1-D_3$, $D_2-D_3$ in $\Phi_3$ are
all zero (\texttt{tests/Hyperelliptic6\_649.m}; the same holds in $\Phi_2\cong\Z/3\times\Z/3\times\Z/90$).

\emph{Conclusion.} Let $i<j$. The class $[D_i-D_j]\in\Jac(Y)(\Q)$ is torsion, of order dividing $6$; its
$3$-part maps injectively to $\Phi_3$ and its image there is zero, so the $3$-part is zero and
$[D_i-D_j]$ lies in $\Jac(Y)(\Q)[2]$, a group of order at most $2$. Three elements $[D_1],[D_2],[D_3]$ of a
coset of a group of order at most two cannot be pairwise distinct, so $D_i\sim D_j$ for some $i<j$, and since
$D_i\neq D_j$ as divisors, $|D_i|$ is a $g^1_2$ on $Y$. \qed

\begin{remark}
The counting argument of \texttt{paper/torsion-counting-hyperelliptic.tex} would need $7$ pairwise related divisors
against the bound $6$; the component group at $3$ reduces the bound to $2$ for these divisors, so three suffice.
The $2$-part cannot be removed the same way: at $3$ the torus has $2$-torsion, and at $2$ the formal group does.
\end{remark}

\end{document}
